\colorlet{line1}{colorBlue}
\colorlet{line2}{colorOrange}
\colorlet{line3}{colorGreen}
\colorlet{line4}{colorGray}

\def\ateachpoint{10}

\pgfplotscreateplotcyclelist{LineCycle}{%
    {line1, thick, mark=x, mark options={solid}, mark size = 2pt, mark repeat=75, filter discard warning=false, unbounded coords=discard},
    {line2, thick, densely dotted, mark=triangle, mark options={solid}, mark size = 1pt, mark repeat=75, filter discard warning=false, unbounded coords=discard},
    {line3, thick, densely dashed, mark=square, mark options={solid}, mark size = 1pt, mark repeat=75, filter discard warning=false, unbounded coords=discard},
    {line4, thick, thin, mark=Mercedes star, mark options={solid}, mark size = 4pt, mark repeat=70, each nth point={\ateachpoint}, filter discard warning=false, unbounded coords=discard},
    {line1, very thick, dashdotted, mark=x, mark options={solid}, mark size = 2pt, mark repeat=45, each nth point={\ateachpoint}, filter discard warning=false, unbounded coords=discard},
    {line2, very thick, dashdotted, mark=x, mark options={solid}, mark size = 2pt, mark repeat=45, each nth point={\ateachpoint}, filter discard warning=false, unbounded coords=discard}
}

\def\ateachpointLoose{30}
\pgfplotscreateplotcyclelist{LineCycleLoose}{%
    {line1, thick, mark=x, mark options={solid}, mark size = 4pt, mark repeat=15, each nth point={\ateachpointLoose}, filter discard warning=false, unbounded coords=discard},
    {line2, thick, densely dotted, mark=triangle, mark options={solid}, mark size = 2pt, mark repeat=17, each nth point={\ateachpointLoose}, filter discard warning=false, unbounded coords=discard},
    {line3, thick, densely dashed, mark=square, mark options={solid}, mark size = 2pt, mark repeat=19, each nth point={\ateachpointLoose}, filter discard warning=false, unbounded coords=discard},
    {line4, thick, thin, mark=Mercedes star, mark options={solid}, mark size = 2pt, mark repeat=16, each nth point={\ateachpointLoose}, filter discard warning=false, unbounded coords=discard},
    {line1, very thick, dashdotted, mark=x, mark options={solid}, mark size = 2pt, mark repeat=18, each nth point={\ateachpointLoose}, filter discard warning=false, unbounded coords=discard},
    {line2, very thick, dashdotted, mark=x, mark options={solid}, mark size = 2pt, mark repeat=15, each nth point={\ateachpointLoose}, filter discard warning=false, unbounded coords=discard}
}

\pgfplotsset{
    /pgfplots/area style/.style={%
        /pgfplots/cycle list={%
            {line1, mark=none, each nth point={\ateachpoint}},%
            {line2, mark=none, each nth point={\ateachpoint}},%
            {line3, mark=none, each nth point={\ateachpoint}},%
            {line4, mark=none, each nth point={\ateachpoint}}%
        }%
    }%
}%

\pgfplotsset{
   boxplot/every median/.style={ultra thick}
}

\pgfplotsset{
    compat=1.18,
    /pgfplots/boxplot legend/.style={
        /pgfplots/legend image code/.code={%
        \draw[##1,/tikz/.cd,yshift=-0.25em]
            (0cm,0cm) rectangle (6pt,6pt);},
    },
}

\pgfplotsset{
    compat=1.18,
	legend image code/.code={
		\draw[mark repeat=2, mark phase=2, rounded corners=2pt,]
		plot coordinates {
			(0cm,0cm)
			(0.15cm,0cm)        %% default is (0.3cm,0cm)
			(0.15cm,0cm)         %% default is (0.6cm,0cm)
		};%
	}
}

\pgfplotsset{
    every axis legend/.style={
        draw=legendBorderColor,
        fill=legendFillColor,
        fill opacity=0.1,
        text opacity=1,
        draw opacity=1, 
        legend columns=1,
        inner sep=2pt,
        /tikz/rounded corners=2pt,
    }
}

\pgfplotsset{
    every axis legend/.append style={
        legend pos=north east,
        legend cell align=left,
    }
}

\pgfplotsset{
    every axis/.style={
        tick align=inside,
        tick pos=both,
        tick style={line width=0.5pt, color=labelColor},
        tick label style={font=\small, color=labelColor},
        label style={font=\small, color=labelColor},
        title style={font=\small},
        legend style={font=\small, text={labelColor}},
        major grid style={dashed, labelColor},
        grid style=dashed,
        xminorgrids,
        yminorgrids,
        enlargelimits=0.05,
        clip=false,
        scaled ticks=false,
        axis line style={labelColor}
    }
}

\tikzstyle{every pin}=[
    labelColor,
    draw=legendBorderColor,
    fill=legendFillColor,
    font=\footnotesize,
    fill opacity=0.1,
    draw opacity=0.4,
    text opacity=1,
    pin distance = 1.6em,
    rounded corners=2pt, 
    inner sep=4pt,
    pin edge={labelColor, dashed, thin}
]

\tikzset{
    plotv/.style={
        domain=\ymin:\ymax,
        samples=2,
        thin,
        dash pattern=on 2pt off 8pt,
        color=labelColor
    },
    ploth/.style={
        domain=\xmin:\xmax,
        samples=2,
        thin,
        dash pattern=on 2pt off 8pt,
        color=labelColor
    },
}

%%%%%%%%%%%%%%%%%%%%%%%% not used yet %%%%%%%%%%%%%%%%%%%%%%%%

% \makeatletter
% \pgfplotsset{
%     boxplot prepared from table/.code={
%         \def\tikz@plot@handler{\pgfplotsplothandlerboxplotprepared}%
%         \pgfplotsset{
%             /pgfplots/boxplot prepared from table/.cd,
%             #1,
%         }
%     },
%     /pgfplots/boxplot prepared from table/.cd,
%         table/.code={\pgfplotstablecopy{#1}\to\boxplot@datatable},
%         row/.initial=0,
%         make style readable from table/.style={
%             #1/.code={
%                 \pgfplotstablegetelem{\pgfkeysvalueof{/pgfplots/boxplot prepared from table/row}}{##1}\of\boxplot@datatable
%                 \pgfplotsset{boxplot/#1/.expand once={\pgfplotsretval}}
%             }
%         },
%         make style readable from table=lower whisker,
%         make style readable from table=upper whisker,
%         make style readable from table=lower quartile,
%         make style readable from table=upper quartile,
%         make style readable from table=median,
%         make style readable from table=lower notch,
%         make style readable from table=upper notch
% }

% \makeatother
%     \pgfmathdeclarefunction{fpumod}{2}{%
%         \pgfmathfloatdivide{#1}{#2}%
%         \pgfmathfloatint{\pgfmathresult}%
%         \pgfmathfloatmultiply{\pgfmathresult}{#2}%
%         \pgfmathfloatsubtract{#1}{\pgfmathresult}%
%         % replaced `0' by `5' to make it work for this problem
%         \pgfmathfloatifapproxequalrel{\pgfmathresult}{#2}{\def\pgfmathresult{5}}{}%
%     }

% % defining the new dimensions and parameters
% \newlength{\hatchspread}
% \newlength{\hatchthickness}
% \newlength{\hatchshift}
% \newcommand{\hatchcolor}{}
% % declaring the keys in tikz
% \tikzset{
%     hatchspread/.code={\setlength{\hatchspread}{#1}},
%     hatchthickness/.code={\setlength{\hatchthickness}{#1}},
%     hatchshift/.code={\setlength{\hatchshift}{#1}},% must be >= 0
%     hatchcolor/.code={\renewcommand{\hatchcolor}{#1}}
% }
% % setting the default values
% \tikzset{
%     hatchspread=3pt,
%     hatchthickness=2pt,
%     hatchshift=0pt,% must be >= 0
%     hatchcolor=white
% }
         
% % declaring the pattern
% \pgfdeclarepatternformonly[\hatchspread,\hatchthickness,\hatchshift,\hatchcolor]% variables
%    {custom north west lines}% name
%    {\pgfqpoint{\dimexpr-2\hatchthickness}{\dimexpr-2\hatchthickness}}% lower left corner
%    {\pgfqpoint{\dimexpr\hatchspread+2\hatchthickness}{\dimexpr\hatchspread+2\hatchthickness}}% upper right corner
%    {\pgfqpoint{\dimexpr\hatchspread}{\dimexpr\hatchspread}}% tile size
%    {% shape description
%     \pgfsetlinewidth{\hatchthickness}
%     \pgfpathmoveto{\pgfqpoint{0pt}{\dimexpr\hatchspread+\hatchshift}}
%     \pgfpathlineto{\pgfqpoint{\dimexpr\hatchspread+0.15pt+\hatchshift}{-0.15pt}}
%     \ifdim \hatchshift > 0pt
%       \pgfpathmoveto{\pgfqpoint{0pt}{\hatchshift}}
%       \pgfpathlineto{\pgfqpoint{\dimexpr0.15pt+\hatchshift}{-0.15pt}}
%     \fi
%     \pgfsetstrokecolor{\hatchcolor}
% %    \pgfsetdash{{1pt}{1pt}}{0pt}% dashing cannot work correctly in all situation this way
%     \pgfusepath{stroke}
%    }

% \pgfdeclarepatternformonly[\hatchspread,\hatchthickness,\hatchshift,\hatchcolor]% variables
%    {custom north east lines}% name
%    {\pgfqpoint{\dimexpr-2\hatchthickness}{\dimexpr-2\hatchthickness}}% lower left corner
%    {\pgfqpoint{\dimexpr\hatchspread+2\hatchthickness}{\dimexpr\hatchspread+2\hatchthickness}}% upper right corner
%    {\pgfqpoint{\dimexpr\hatchspread}{\dimexpr\hatchspread}}% tile size
%    {% shape description
%     \pgfsetlinewidth{\hatchthickness}
%     \pgfpathmoveto{\pgfqpoint{\dimexpr\hatchshift-0.15pt}{-0.15pt}}
%     \pgfpathlineto{\pgfqpoint{\dimexpr\hatchspread+0.15pt}{\dimexpr\hatchspread-\hatchshift+0.15pt}}
%     \ifdim \hatchshift > 0pt
%       \pgfpathmoveto{\pgfqpoint{-0.15pt}{\dimexpr\hatchspread-\hatchshift-0.15pt}}
%       \pgfpathlineto{\pgfqpoint{\dimexpr\hatchshift+0.15pt}{\dimexpr\hatchspread+0.15pt}}
%     \fi
%     \pgfsetstrokecolor{\hatchcolor}
% %    \pgfsetdash{{1pt}{1pt}}{0pt}% dashing cannot work correctly in all situation this way
%     \pgfusepath{stroke}
%    }

% \pgfdeclarepatternformonly{mydots}
%     {\pgfqpoint{-1pt}{-1pt}}
%     {\pgfqpoint{1pt}{1pt}}
%     {\pgfqpoint{2pt}{2pt}}% original definition: \pgfqpoint{3pt}{3pt}
% {%
%   \pgfpathcircle{\pgfqpoint{0pt}{0pt}}{.8pt}%
%   \pgfusepath{fill}%
% }%

% % declaring the pattern
% \pgfdeclarepatternformonly%[\hatchspread,\hatchthickness,\hatchshift,\hatchcolor]% variables
%    {mynorthwestlines}% name
%    {\pgfqpoint{-4pt}{-4pt}}% lower left corner
%    {\pgfqpoint{8pt}{8pt}}% upper right corner
%    {\pgfqpoint{4pt}{4pt}}% tile size
%    {% shape description
%     \pgfsetlinewidth{2pt}
%     \pgfpathmoveto{\pgfqpoint{0pt}{4pt}}
%     \pgfpathlineto{\pgfqpoint{4.15pt}{-0.15pt}}
%     \pgfsetstrokecolor{white}
% %    \pgfsetdash{{1pt}{1pt}}{0pt}% dashing cannot work correctly in all situation this way
%     \pgfusepath{stroke}
%    }

% Style to select only points from #1 to #2 (inclusive)
% \pgfplotsset{select coords between index/.style 2 args={
%     x filter/.code={
%         \ifnum\coordindex<#1\def\pgfmathresult{}\fi
%         \ifnum\coordindex>#2\def\pgfmathresult{}\fi
%     }
% }}

% \newsavebox{\measuredSize}
% \newcommand{\resizeToWidth}[2]{%
%     \pgfmathsetmacro{\pgfplotswidth}{#2}%
%     \begin{lrbox}{\measuredSize}#1\end{lrbox}%
%     \pgfmathsetmacro{\pgfplotswidth}{2*\pgfplotswidth-\wd\measuredSize}%
%     #1%
% }

% \tikzset{every picture/.style={scale=1.4, transform shape}}

% \pgfplotsset{
%     discard if/.style 2 args={
%         x filter/.code={
%             \edef\tempa{\thisrow{#1}}
%             \edef\tempb{#2}
%             \ifx\tempa\tempb
%                 \def\pgfmathresult{inf}
%             \fi
%         }
%     },
%     discard if not/.style 2 args={
%         x filter/.code={
%             \edef\tempa{\thisrow{#1}}
%             \edef\tempb{#2}
%             \ifx\tempa\tempb
%             \else
%                 \def\pgfmathresult{inf}
%             \fi
%         }
%     }
% }

% \pgfplotscreateplotcyclelist{AreaCycle}{%
%     {line1, mark=none, pattern=crosshatch, pattern color = colorGray, each nth point={\ateachpoint}},
%     {line2, mark=none, pattern=mydots, pattern color = colorGreen, each nth point={\ateachpoint}},
%     {line3, mark=none, each nth point={\ateachpoint}},
%     {line4, mark=none, each nth point={\ateachpoint}}
% }
